% !TeX root = ../main.tex

\chapter{总结与展望}

四种nPDF的的核修正因子存在显著差异，重点研究这些差异如何体现在双喷注的观测量上。我们希望我们的研究有助于未来的实验测量，以阐明这些核效应的范围。



	原子核部分子分布相对于自由核子分布的偏离，源于将核子结合在一起的额外非微扰动力学，通常可以用修正因子$r_i^A(x,Q^2)={f_i^{A,p}(x,Q^2)}/{f_i^p(x,Q^2)}$来定量描述，其中$A$为原子核质量数，$i$表示部分子类型，$x$表示动量分数，$Q^2$表示分辨标度。
	从实验测量中提取$r_i^A(x,Q^2)$是获得核修正的一种不可替代的方式，而且这依赖于微扰量子色动力学中的因子化。在实际的可观测量（例如微分散射截面）中，不同$x$区域、不同探测标度和不同部分子类型的核效应通常会混合在一起。然而，一个能够在领头阶级别上准确解析$x$和$Q^2$的可观测量，就会提供$r_i^A(x,Q^2)$对$x$和$Q^2$依赖性的更有效和详细的动力学扫描。

	在未来的高精度下，LHC实验对双喷注的三重微分散射截面$R_{p\textrm{Pb}}$ 的测量有望在各种nPDF参数化的集合上提供显著且多维度的约束。
	更重要的是，这些测量将为确认不同$x$区域的各种核效应（如遮蔽效应、反遮蔽效应、EMC效应和费米运动效应）以及它们随着探测标度$Q^2$的变化提供有效的方法。
	值得注意的是，所研究的三重微分散射截面可以直接推广和应用于其它过程的研究，例如矢量玻色子标记的喷注产生~\cite{Ma:2018tjv,Zhang:2018urd,Zhang:2021oki}和重夸克双喷注的产生~\cite{Dai:2018mhw}。
	我们还希望提供一个对初态冷核物质效应有更准确描述的基准，以便更好地理解相对论重离子碰撞中末态的喷注淬火现象。

	本章的研究中，我们关注了LHC实验上$p$Pb碰撞中双喷注的产生过程作为原子核夸克和胶子分布的探针。为了很好地解析初态部分子的动量分数以及探测标度，我们研究了$pp$和$p$Pb碰撞中几种类型的三重微分截面，包括以双喷注变量$V^{(3)}=\{p_{T,\textrm{avg}},y_b,y^\}$，$\{X_B,X_A,y^\}$和$\{X_B,E_T^{{JJ}},p_{T,\textrm{avg}}\}$为分区间的散射截面。这些截面在微扰量子色动力学中被计算到次领头阶。我们在$p$Pb的计算中使用了四组原子核部分子分布函数（nPDFs），分别是EPPS16、nCTEQ15、TUJU19和nIMParton16。特别说明，可观测的核修正因子$R_{p\textrm{Pb}}$的三重微分截面，尤其是作为$X_B$变量的函数，能够很好地展示$x$值从小到大的nPDF修正因子$r_i^{A}(x,Q^2)$的图像。基于此，我们便可以很好地解释不同nPDFs预测的$R_{p\textrm{Pb}}$之间的差异。特别是，通过选取$V^{(3)}=\{X_B,E_T^{{JJ}},p_{T,\textrm{avg}}\}$，$R_{p\textrm{Pb}}$随着$p_{T,\textrm{avg}}$的变化甚至能够反映$r_i^{A}(x,Q^2)$的微妙的标度变化。


